\documentclass[11pt]{article}
\usepackage[margin=1in]{geometry}
\usepackage[T1]{fontenc}
\usepackage[utf8]{inputenc}
\usepackage{amsmath,amssymb,amsthm}
\usepackage{graphicx}
\usepackage{xcolor}
\usepackage{microtype}
\usepackage[colorlinks=true,linkcolor=blue!55!black,citecolor=blue!55!black,urlcolor=blue!55!black]{hyperref}

\newtheorem{theorem}{Theorem}
\newtheorem{lemma}{Lemma}
\newtheorem{corollary}{Corollary}
\newcommand{\one}{\mathbf 1}

\title{Continuous boundary data do not rescue interval square roots in $C(K)$}
\author{Candidate counterexample packet for arXiv:1604.01211}
\date{11 August 2026}

\begin{document}
\maketitle

\begin{center}
\fbox{\parbox{0.92\textwidth}{\textbf{Status: candidate full counterexample, likely valid, subject to human review.}
For every compact Hausdorff space containing a nontrivial convergent sequence,
there are continuous $f,g$ satisfying the source paper's necessary pointwise
conditions $g\geq0$ and $|f|\leq g$, but $[f,g]$ has no square root which is an
interval.  Consequently, among first-countable compacta, finiteness is exactly
the condition under which those necessary conditions suffice for all continuous
boundary data.}}
\end{center}

\section{The source question and the scoped answer}

For a compact Hausdorff space $K$, an interval in the real Banach algebra
$C(K)$ is
\[
 [f,g]=\{h\in C(K):f\leq h\leq g\}.
\]
Moreno and Schneider call an interval $I$ a square root of $[f,g]$ when
$I\cdot I=[f,g]$, with pointwise multiplication.  In Section~5 they ask:
\begin{quote}
``Which intervals have a square root that is an interval?''
\end{quote}
They prove that $g\geq0$ and $|f|\leq g$ are necessary, that $f\geq0$ is
sufficient, and that $-g\leq f\leq0<g$ is sufficient when $g$ is continuous
\cite[Proposition~5.1]{source}.  The result below shows that the necessary
conditions are not sufficient even when \emph{both} boundary functions are
continuous.  It also gives a complete answer to the corresponding uniform
sufficiency question for first-countable compacta.

Figure~\ref{fig:source} reproduces the source question and the neighboring
necessary and sufficient conditions from printed page~14.

\begin{figure}[ht]
\centering
\includegraphics[width=0.94\textwidth]{figures/open_problem_crop.png}
\caption{The square-root question and Proposition~5.1 in the source paper.}
\label{fig:source}
\end{figure}

\clearpage
\section{Idea of the proof}

Take $g\equiv1$ and make a continuous lower boundary $f$ alternate between
small positive and small negative values along a sequence converging to
$x_0$.  If an interval $I$ squared to $[f,\one]$, then the member $f$ itself
would factor as $f=hk$ with $h,k\in I$.

At a point where $f=a>0$, the scalar interval $I(x)$ cannot cross zero, and
attaining the lower endpoint $a$ forces both factors to have modulus
$\sqrt a$.  Hence both $h$ and $k$ tend to zero along the positive
subsequence.  At a point where $f=a<0$, attaining the lower endpoint of the
scalar square forces one factor to be an endpoint of modulus $1$.  Along the
negative subsequence, therefore, one of $h,k$ always has modulus $1$.  These
two conclusions are incompatible with continuity at $x_0$.

\section{The scalar endpoint obstruction}

\begin{lemma}\label{lem:scalar}
Let $-1\leq a\leq1$, $a\neq0$, and let $J=[c,d]$ be a nonempty compact real
interval such that
\[
 J\cdot J=[a,1].
\]
If $u,v\in J$ and $uv=a$, then:
\begin{enumerate}
\item if $a>0$, one has $|u|=|v|=\sqrt a$;
\item if $a<0$, one has $\max\{|u|,|v|\}=1$.
\end{enumerate}
\end{lemma}

\begin{proof}
Suppose first that $a>0$.  Since $J\cdot J$ contains neither zero nor a
negative number, $J$ lies strictly on one side of zero.  If $0<c\leq d$,
then $J\cdot J=[c^2,d^2]$, so $c^2=a$, $d^2=1$, and $uv=a$ forces
$u=v=c$.  If $c\leq d<0$, then $J\cdot J=[d^2,c^2]$, and the symmetric
argument gives $u=v=d$.  This proves the first assertion.

If $a<0$, then $c<0<d$ and the lower endpoint of $J\cdot J$ is $cd=a$,
whereas its upper endpoint is $\max\{c^2,d^2\}=1$.  For $uv=a$, the two
factors have opposite signs.  Assuming $u<0<v$, we have
$|u|\leq|c|$ and $v\leq d$; equality in
$|u|v=|a|=|c|d$ forces $u=c$ and $v=d$.  The other ordering is identical.
Thus $\max\{|u|,|v|\}=\max\{|c|,d\}=1$.
\end{proof}

\section{Continuous counterexamples}

\begin{theorem}\label{thm:main}
Let $K$ be a compact Hausdorff space containing a sequence of distinct points
$(x_n)$ converging to a point $x_0\notin\{x_n:n\geq1\}$.  There exists a
continuous function $f:K\to[-1,1]$ such that
\[
 \one\geq |f|,
\]
but the continuous interval $[f,\one]$ has no square root which is an
interval in $C(K)$.
\end{theorem}

\begin{proof}
The set $S=\{x_0\}\cup\{x_n:n\geq1\}$ is closed in $K$.  Define
\[
 r(x_0)=0,\qquad r(x_n)=\frac{(-1)^n}{n+1}.
\]
This is a continuous function $S\to[-1,1]$.  Compact Hausdorff spaces are
normal, so the Tietze extension theorem gives an
$f\in C(K,[-1,1])$ extending $r$ \cite[Section~33]{munkres}.

Assume for contradiction that an interval $I$ satisfies
\[
 I\cdot I=[f,\one].
\]
Replace any presentation of $I$ by its canonical presentation.  By
\cite[Lemma~3.1(i)]{source}, evaluation at each $x\in K$ gives a scalar
interval $I(x)$ satisfying
\[
 I(x)\cdot I(x)=[f(x),1].
\]
Since $f\in[f,\one]=I\cdot I$, there are $h,k\in I$ with $hk=f$.

For even $n$, $f(x_n)>0$.  Lemma~\ref{lem:scalar} gives
\[
 |h(x_n)|=|k(x_n)|=\sqrt{f(x_n)}\longrightarrow0.
\]
Continuity along the even subsequence implies
$h(x_0)=k(x_0)=0$.  For odd $n$, however, $f(x_n)<0$, and the same lemma
gives
\[
 \max\{|h(x_n)|,|k(x_n)|\}=1.
\]
Continuity of $h$ and $k$ at $x_0$ makes the left-hand side tend to zero
along the odd subsequence, a contradiction.
\end{proof}

The explicit metrizable example is $K=[-1,1]$, $f(t)=t$, and $g(t)=1$.
Thus no extension theorem is needed for the simplest instance.

\begin{corollary}\label{cor:characterization}
For a first-countable compact Hausdorff space $K$, the following are
equivalent:
\begin{enumerate}
\item $K$ is finite;
\item every continuous interval $[f,g]$ satisfying
$g\geq0$ and $|f|\leq g$ has a square root which is an interval.
\end{enumerate}
\end{corollary}

\begin{proof}
If $K$ is infinite, compactness supplies an accumulation point.  A countable
local base at that point lets one choose a sequence of distinct points
converging to it, so Theorem~\ref{thm:main} disproves (2).

If $K$ is finite, $C(K)$ is a finite Cartesian product of copies of
$\mathbb R$.  For each $x\in K$, the scalar condition
$g(x)\geq0$, $|f(x)|\leq g(x)$ produces a scalar interval whose square is
$[f(x),g(x)]$ (the elementary classification recorded in the proof of
\cite[Proposition~5.1]{source}).  Choosing one such interval independently
at each of the finitely many points gives an interval in $C(K)$ whose square
is $[f,g]$.
\end{proof}

\section{Scope, novelty check, and review recommendation}

The theorem does not classify individual square-rootable intervals on an
arbitrary compactum.  It decisively rules out the most direct completion of
the source proposition---sufficiency of its pointwise necessary inequalities,
even for continuous boundary data---and Corollary~\ref{cor:characterization}
fully characterizes the first-countable compacta on which that uniform
sufficiency statement holds.

On 11 August 2026, exact arXiv-id/title searches in the run registry and all
cheap run indexes found no prior packet.  Searches of the locally parsed arXiv
corpus for the source title, ``square roots of intervals,'' and the necessary
conditions found only the source.  Bounded web searches for the exact question,
the source title with ``square root,'' and close $C(K)$ variants likewise found
the 2016 paper but no later resolution or this continuous-boundary obstruction.
This is a bounded novelty check, not an exhaustive priority claim.

Human review should focus on the use of source Lemma~3.1(i) to identify every
pointwise scalar square and on the equality cases in Lemma~\ref{lem:scalar}.
The explicit $[-1,1]$ example isolates the entire analytic obstruction and
can be checked independently of the Tietze-extension strengthening.

\begin{thebibliography}{9}
\bibitem{source}
J. P. Moreno and R. Schneider,
\emph{Multiplication of convex sets in $C(K)$ spaces},
Studia Math. 232 (2016), 173--187; arXiv:1604.01211.

\bibitem{munkres}
J. R. Munkres,
\emph{Topology}, 2nd ed., Prentice Hall, 2000.
\end{thebibliography}

\end{document}
